\documentclass[a4paper,10pt]{book}
\usepackage[utf8]{inputenc}

\usepackage{listings}
\usepackage{xcolor}

% Configuração do ambiente listings
\lstdefinestyle{mystyle}{
  basicstyle=\ttfamily,
  language=[LaTeX]TeX,
  breaklines=true,
  columns=fullflexible,
  frame=none,%L,%none,%single
  numbers=none, 
  backgroundcolor=\color{gray!10},
  %numbers=left,
  %numberstyle=\tiny,
  keywordstyle=\color{green!50!black},
  commentstyle=\color{magenta}, 
  %stringstyle=\color{orange},
  %identifierstyle=\color{blue},
  %xleftmargin=2ex,
  %xrightmargin=2ex,
  classoffset=1,
  morekeywords={NewEnvBoxFormalTabGlobal,DefCtxVar,SetLenCtxVar},
  keywordstyle=\color{red!50!black},
  classoffset=2,
  morekeywords={lipsum,tcblistof,ding},
  keywordstyle=\color{blue!50!black},
  classoffset=0,
}


\usepackage[
TitleColor=blue!50!black,
TitleFontSize=30,
PreTitleColor=black,
PreTitleFontSize=90,
BarWidth=2cm,
BarTopColor=blue!20,
BarBottomColor=red!20,
BarTextColor=white,
BarTextFontSize=30,
BarTextOffset=1cm,
AfterSpace=0cm
]{chapter-format-rightbar-text}
\usepackage{lipsum}
\usepackage{comment}
\usepackage[margin=25mm]{geometry}
\usepackage{graphicx}
\usepackage{mdframed}
\usepackage{wrapfig}

%opening
\title{chapter-format-rightbar-text package}
\author{Fernando Pujaico Rivera}


\hyphenpenalty=10000
\exhyphenpenalty=10000

\begin{document}

\maketitle

%-------------------------------------------------------------------------------
%   TABLE OF CONTENTS
%-------------------------------------------------------------------------------
\clearpage
\pagestyle{empty} % No headers %\pagestyle{fancy} % Print headers again
\tableofcontents


%%%%%%%%%%%%%%%%%%%%%%%%%%%%%%%%%%%%%%%%%%%%%%%%%%%%%%%%%%%%%%%%%%%%%%%%%%%%%%%%%%%%%%%%%%%%%%%%%
%%%%%%%%%%%%%%%%%%%%%%%%%%%%%%%%%%%%%%%%%%%%%%%%%%%%%%%%%%%%%%%%%%%%%%%%%%%%%%%%%%%%%%%%%%%%%%%%%
\chapter{Install information}\label{ch:install}

%%%%%%%%%%%%%%%%%%%%%%%%%%%%%%%%%%%%%%%%%%%%%%%%%%%%%%%%%%%%%%%%%%%%%%%%%%%%%%%%%%%%%%%%%%%%%%%%%
\section{Install package}\label{sec:install}
Put the chapter-format-rightbar-text.sty file in any of these locations
\begin{itemize}
\item Put the \texttt{chapter-format-rightbar-text.sty} file in the same path of main tex file, or.
\item Execute the commmand:
\begin{lstlisting}[style=mystyle]
kpsewhich -var-value=TEXMFHOME
\end{lstlisting}
and this returns the path of local tex files. By example, if returns 
\begin{lstlisting}[style=mystyle]
/home/username/texmf
\end{lstlisting}
then, put the \texttt{chapter-format-rightbar-text.sty} file in the directory.
\begin{lstlisting}[style=mystyle]
/home/username/texmf/tex/latex/chapter-format-rightbar-text/chapter-format-rightbar-text.sty
\end{lstlisting}
\end{itemize}
%%%%%%%%%%%%%%%%%%%%%%%%%%%%%%%%%%%%%%%%%%%%%%%%%%%%%%%%%%%%%%%%%%%%%%%%%%%%%%%%%%%%%%%%%%%%%%%%%
\subsection{Install and compile when we use a custom path}
If we put the \texttt{chapter-format-rightbar-text.sty} style file in the path \texttt{/path/of/styles}.
Then, we need follow the next steps to compile the \texttt{maincode.tex} main tex file that use the\\% 
\noindent\verb|\usepackage{chapter-format-rightbar-text}| code.
\begin{lstlisting}[style=mystyle]
export TEXINPUTS=/path/of/styles:$TEXINPUTS
pdflatex maincode
pdflatex maincode
pdflatex maincode
\end{lstlisting}
The last commands are executed from the path of \texttt{maincode.tex} file.


%%%%%%%%%%%%%%%%%%%%%%%%%%%%%%%%%%%%%%%%%%%%%%%%%%%%%%%%%%%%%%%%%%%%%%%%%%%%%%%%%%%%%%%%%%%%%%%%%
\section{Load the package}
To load the package use the next command in the preamble of main tex document.
\begin{lstlisting}[style=mystyle]
\usepackage{chapter-format-rightbar-text}
\end{lstlisting}

The command \texttt{usepackage} find the \texttt{chapter-format-rightbar-text.sty} file 
in the directories listed in the Section \ref{sec:install}.
By other side, if we locate the \texttt{chapter-format-rightbar-text.sty} file in
\texttt{/some/path}, then we can load the package using the next command.
\begin{lstlisting}[style=mystyle]
\usepackage{/some/path/chapter-format-rightbar-text}
\end{lstlisting}



%%%%%%%%%%%%%%%%%%%%%%%%%%%%%%%%%%%%%%%%%%%%%%%%%%%%%%%%%%%%%%%%%%%%%%%%%%%%%%%%%%%%%%%%%%%%%%%%%
\section{Using the package}
They are shown some examples use of all macros in the package.
All the code should setted before the \verb|\begin{document}|.
The next code generates a chapter format like the used in this document.

\begin{lstlisting}[style=mystyle]
\usepackage[
TitleColor=blue!50!black,
TitleFontSize=30,
PreTitleColor=black,
PreTitleFontSize=90,
BarWidth=2cm,
BarTopColor=blue!20,
BarBottomColor=red!20,
BarTextColor=white,
BarTextFontSize=30,
BarTextOffset=1cm,
AfterSpace=0cm
]
{chapter-format-rightbar-text}
\end{lstlisting}

%%
In the next pages will be showed the results of the next code.
\begin{lstlisting}[style=mystyle]
\SetChapterFormatRightBarText{Some example text with a very long size}
\chapter{Title example of chapter with number}
The \lipsum[1][1-4]
\section{First title section}
The \lipsum[1-2]
\section{Second title section}
The \lipsum[1-2]

\SetChapterFormatRightBarText{Perfect}
\chapter*{Title example of chapter without number}
\addcontentsline{toc}{chapter}{Title example of chapter without number}
The \lipsum[1][1-4]
\section*{Third title section}
\addcontentsline{toc}{section}{Third title section}
The \lipsum[1-2]
\section*{Fourth title section}
\addcontentsline{toc}{section}{Fourth title section}
The \lipsum[1-2]
\end{lstlisting}

%%%%%%%%%%%%%%%%%%%%%%%%%%%%%%%%%%%%%%%%%%%%%%%%%%%%%%%%%%%%%%%%%%%%%%%%%%%%%%%%%%%%%%%%%%%%%%%%%
%%%%%%%%%%%%%%%%%%%%%%%%%%%%%%%%%%%%%%%%%%%%%%%%%%%%%%%%%%%%%%%%%%%%%%%%%%%%%%%%%%%%%%%%%%%%%%%%%
\SetChapterFormatRightBarText{Some example text with a very long size}
\chapter{Title example of chapter with number}
The \lipsum[1][1-4]
%%%%%%%%%%%%%%%%%%%%%%%%%%%%%%%%%%%%%%%%%%%%%%%%%%%%%%%%%%%%%%%%%%%%%%%%%%%%%%%%%%%%%%%%%%%%%%%%%
\section{First title section}
The \lipsum[1-2]
%%%%%%%%%%%%%%%%%%%%%%%%%%%%%%%%%%%%%%%%%%%%%%%%%%%%%%%%%%%%%%%%%%%%%%%%%%%%%%%%%%%%%%%%%%%%%%%%%
\section{Second title section}
The \lipsum[1-2]


%%%%%%%%%%%%%%%%%%%%%%%%%%%%%%%%%%%%%%%%%%%%%%%%%%%%%%%%%%%%%%%%%%%%%%%%%%%%%%%%%%%%%%%%%%%%%%%%%
%%%%%%%%%%%%%%%%%%%%%%%%%%%%%%%%%%%%%%%%%%%%%%%%%%%%%%%%%%%%%%%%%%%%%%%%%%%%%%%%%%%%%%%%%%%%%%%%%
\SetChapterFormatRightBarText{Perfect}
\chapter*{Title example of chapter without number}
\addcontentsline{toc}{chapter}{Title example of chapter without number}
The \lipsum[1][1-4]
%%%%%%%%%%%%%%%%%%%%%%%%%%%%%%%%%%%%%%%%%%%%%%%%%%%%%%%%%%%%%%%%%%%%%%%%%%%%%%%%%%%%%%%%%%%%%%%%%
\section*{Third title section}
\addcontentsline{toc}{section}{Third title section}
The \lipsum[1-2]
%%%%%%%%%%%%%%%%%%%%%%%%%%%%%%%%%%%%%%%%%%%%%%%%%%%%%%%%%%%%%%%%%%%%%%%%%%%%%%%%%%%%%%%%%%%%%%%%%
\section*{Fourth title section}
\addcontentsline{toc}{section}{Fourth title section}
The \lipsum[1-2]

\end{document}
